\documentclass[11pt]{article}
\usepackage[T1]{fontenc}
\usepackage[utf8]{inputenc}
\usepackage[english]{babel}
\usepackage[margin=1in]{geometry}
\usepackage{amsmath,amssymb}
\usepackage{enumitem}
\usepackage{microtype}

\ifdefined\pdfcatalog
  \pdfcatalog{/Lang (en-US)}
\fi

\setlength{\parindent}{0pt}
\setlength{\parskip}{0.55em}
\setlength{\emergencystretch}{2em}
\setlist{nosep,leftmargin=*}

\newcommand{\findinglabel}[1]{\textbf{#1}}
\newcommand{\classification}[1]{\textit{Classification: #1.}}

\title{Technical Review of ``Early-Stopping Activation Steering for Factual Preference Alignment in Vietnamese Medical Language Models''}
\author{Manuscript review}
\date{August 27, 2026}

\begin{document}
\selectlanguage{english}
\maketitle

\section*{Overall assessment}

The revision makes useful local changes: the RAG prompt-token comparison now uses a 35.0-token instruction-only baseline, the natural-completion discussion reports condition-specific EOS rates, and the bounded 200-token McNemar calculation remains internally consistent.  The manuscript also attempts to add multiplicity-adjusted inference for natural completion.

The paper is nevertheless not submission-ready.  The newly reported placebo distribution is impossible as written and contradicts its own table; the claimed empirical randomization value cannot be obtained from the stated 20-direction procedure.  The new natural-completion significance table appears to reuse the bounded-generation confidence interval and McNemar result even though the displayed effect is different.  The implementation target is Layer~8 in most sections but $l=16$ in the steering equation.  In addition, a clean build fails because a \texttt{table} environment is closed as \texttt{table*}, and an unescaped ampersand in another caption produces a second LaTeX error.  Mechanistic, clinical-validity, RAG, timing, dataset, and reproducibility limitations remain.  Major revision is required.

\section*{Major technical findings}

\subsection*{1. The placebo distribution is mathematically impossible and internally contradictory}

\classification{Definite quantitative and data-reporting errors}

\findinglabel{Location.} Abstract; contribution~4; Placebo Control Distribution \& Real BM25 RAG Results; Table~\texttt{tab:baselines}; Conclusion.

\findinglabel{Problem.} The Abstract, Results, and Conclusion report a random-vector mean of $32.86\%\pm0.24\%$, while the Results simultaneously states that the 20 values range from 45.00\% to 63.00\%.  A mean of 32.86\% cannot lie below every observation in that range.  Table~\texttt{tab:baselines} instead reports $56.50\%\pm5.11\%$.  These alternatives imply radically different standardized distances:
\[
\frac{77.20-32.86}{0.24}=184.75,
\qquad
\frac{77.20-56.50}{5.11}\approx4.05,
\]
neither of which equals the stated $185.77\sigma$ using the displayed rounded quantities.  The meaning of the ``$\pm$'' term is also not defined as a standard deviation, standard error, or confidence interval.

The claimed empirical randomization $p=0.0099$ is incompatible with the stated comparison against only $R=20$ sampled directions if the usual plus-one correction is used.  With no random direction exceeding the target, the smallest attainable value is
\[
p_{\min}=\frac{1}{R+1}=\frac{1}{21}\approx0.0476.
\]
A value of 0.0099 would require a different resampling experiment that is not described.

\findinglabel{Why it matters.} Directional specificity is a headline contribution, but the current summary does not identify the actual random-direction results or a valid inferential procedure.

\findinglabel{Suggested fix.} Regenerate the placebo section from the 20 per-direction outputs.  Publish every seed and direction-level metric; report the mean, sample standard deviation, range, and uncertainty with explicit definitions.  Use a prespecified randomization test whose resolution matches the number of sampled directions, or run substantially more directions.  Remove all $\sigma$ and $p$ claims until one authoritative distribution and test are established.

\subsection*{2. The natural-completion significance table mixes results from two different experiments}

\classification{Definite statistical reporting and provenance error}

\findinglabel{Location.} bounded-generation significance paragraph; Table~\texttt{tab:mcnemar}; Natural Termination Validation; Abstract and Conclusion.

\findinglabel{Problem.} For bounded generation, the reported discordant counts $b=16$ and $c=37$ imply
\[
\frac{c-b}{500}=\frac{21}{500}=4.20\text{ pp}
\]
and produce the stated exact McNemar value $p\approx0.0055$.  Table~\texttt{tab:mcnemar} is instead captioned as natural long-sequence completion and displays Hard Cutoff changing 327/500 to 353/500, a 26-record or 5.20-point effect.  Nevertheless, its Hard Cutoff row reuses $p=0.0055$ and the interval $[+1.37,+7.03]$~pp from the bounded 4.20-point analysis.  That interval is centered at 4.20 points rather than the displayed 5.20 points.  The table supplies no natural-completion discordant counts from which its McNemar values can be checked.

\findinglabel{Why it matters.} This changes the primary effect estimate, test population, schedule, generation cap, and multiplicity-adjusted conclusion.  The claimed significant natural-completion result is not supported by an identifiable paired dataset in the manuscript.

\findinglabel{Suggested fix.} Separate bounded and natural-completion inference into independently generated tables.  For each natural schedule, report the complete $2\times2$ paired transition table, raw counts, paired effect and CI, exact test call, and Holm procedure.  Generate captions, intervals, and prose from the same versioned analysis artifact and remove copied quantities.

\subsection*{3. The intervention layer and activation symbols are contradictory}

\classification{Definite implementation contradiction and notation drift}

\findinglabel{Location.} Abstract; vector controls; Early-Stopping Steering Hook \& Saturation Diagnostics; Experimental Setup; activation Results; layer sweep.

\findinglabel{Problem.} The paper repeatedly states that extraction, steering, and diagnostics use Layer~8.  The implementation-facing hook definition instead says that residual activations are modified at target layer $l=16$.  The source does not state whether layer numbering is zero-based or one-based, so this cannot be reconciled as an indexing convention.

There is also required descriptor drift in $h$.  Equation~1 defines $h_l(x_i,y_i)$ as the final-token layer output of a completed prompt--response sequence.  The hook section later uses $h_l^{(t)}$ for the pre-intervention residual stream and $\hat h_l^{(t)}$ for the post-intervention value.  The Abstract calls $\|h_8^{(t)}\|_2$ a post-hook norm even though the hat is omitted, leaving the diagnostic's measurement point ambiguous.

\findinglabel{Why it matters.} Layer~8 and Layer~16 are different interventions, and pre-hook versus post-hook norms answer different questions.  Independent reproduction and interpretation of the mechanism are impossible without the exact tensor and indexing convention.

\findinglabel{Suggested fix.} Identify the exact Qwen module path, layer index used by code, human-readable layer number, hook input/output tensor, token positions, and pre/post-hook logging order.  Use $h_l^{(t)}$ only for the unmodified value and $\hat h_l^{(t)}$ for the modified value throughout.  Correct either $l=16$ or every Layer~8 claim from the actual run configuration.

\subsection*{4. Activation norms do not demonstrate saturation prevention or stabilization}

\classification{Unsupported mechanism claim}

\findinglabel{Location.} Abstract; Introduction; contribution~3; activation diagnostics and Discussion.

\findinglabel{Problem.} A scalar $\ell_2$ norm cannot establish representation saturation, departure from a high-dimensional activation manifold, or recovery of representation dynamics.  At $t=50$, continuous steering is 1.97 units above the 54.21 baseline, whereas linear decay is 3.01 units below it.  Thus the linear-decay value is farther from the displayed baseline than the continuous value is; calling 51.20 ``converging back'' or ``stabilization'' requires a prespecified acceptable band or functional criterion.  Only selected values at $t=10$ and $t=50$ are reported, and the hard-cutoff change from 56.28 to 52.75 has no time indices.

For $t>K$, both early-stop schedules have zero direct intervention.  Their later free-generation states nevertheless depend on different previously generated tokens, so between-condition norm differences conflate the intervention with divergent token histories.  No dispersion, confidence interval, per-step survivor count, or association with output errors is provided.

\findinglabel{Why it matters.} The central novelty is justified as prevention of an internal failure, but the current diagnostic establishes only under-specified norm differences after different generation trajectories.

\findinglabel{Suggested fix.} Define saturation operationally before analysis.  Add fixed-token replay or teacher forcing, full pre/post-hook trajectories, uncertainty bands, survivor counts, projection onto $v_{\text{steer}}$, cosine drift, covariance or distributional distance, and response to an additional controlled perturbation.  Relate these measurements to repetition and factual errors.  Until then, use ``norm elevation'' and ``norm undershoot'' rather than ``saturation,'' ``manifold departure,'' or ``stabilization.''

\subsection*{5. Natural completion shows a trade-off and does not establish the superiority of linear decay}

\classification{Unsupported comparative interpretation}

\findinglabel{Location.} Abstract; contribution~2; natural-completion Results; Table~\texttt{tab:long\_gen}; Discussion; Conclusion.

\findinglabel{Problem.} Every steering schedule increases RefPref but lowers reference BERTScore and increases Rep-4 relative to baseline.  Hard Cutoff raises Rep-4 from 4.10\% to 5.35\%, a 30.5\% relative increase; Linear Decay reaches the highest repetition value, 5.37\%.  Baseline is best for BERTScore and Rep-4, while Hard Cutoff is best for RefPref.  The manuscript still calls the repetition change ``minor'' and concludes that linear decay yields consistent improvements.  In the manuscript's own new inference table, Linear Decay versus baseline is not significant ($p=0.1250$), although the table itself must first be corrected as noted above.

The universal 100.00\% EOS statement remains in contribution~2 and Experimental Setup despite the condition-specific 99.80\% values for Continuous Steering and Linear Decay.

\findinglabel{Why it matters.} The data do not show uniform output-quality improvement or establish Linear Decay as the best natural-completion schedule.  The primary method and primary endpoint remain insufficiently separated from exploratory schedule selection.

\findinglabel{Suggested fix.} Prespecify the primary endpoint and schedule using validation data, evaluate that frozen choice once on test data, and report paired CIs for RefPref, BERTScore, and Rep-4.  If claiming acceptable trade-offs, define noninferiority margins for semantic similarity and repetition.  Replace universal EOS and general ``improvement'' language with condition- and metric-specific statements.

\subsection*{6. RefPref remains a semantic proxy rather than clinical factual correctness}

\classification{Unsupported interpretation of a disclosed proxy}

\findinglabel{Location.} title; Abstract; Introduction; contributions; Results; Conclusion; keywords.

\findinglabel{Problem.} Methods appropriately discloses that RefPref is not clinician-adjudicated correctness.  Elsewhere, however, the direction is called ``truthfulness,'' RefPref is repeatedly called ``accuracy'' or ``factual alignment,'' and the work is framed as hallucination mitigation and clinical decision support.  Relative BERTScore similarity to $y^+$ rather than one synthetic $y^-$ cannot identify a wrong dose, unit, negation, contraindication, or interaction mechanism.  The natural-completion results also show that RefPref can increase while BERTScore to the reference decreases.

\findinglabel{Why it matters.} Clinical safety, harmful-error reduction, and deployment suitability cannot be inferred from the current endpoint.

\findinglabel{Suggested fix.} Add blinded clinician-adjudicated correctness and harmful-error rates, stratified by pregnancy, interaction, and dosage cases.  Report reviewer qualifications, agreement, adjudication, and validation of RefPref against these labels.  Otherwise use ``reference-preference rate'' consistently and narrow the title, Abstract, keywords, and conclusions.

\subsection*{7. The RAG comparison remains under-specified and statistically unsupported}

\classification{Likely weak baseline and unsupported causal/efficiency claims}

\findinglabel{Location.} Abstract; Experimental Setup; RAG Results; Table~\texttt{tab:baselines}; Conclusion.

\findinglabel{Problem.} The prompt-token comparison is now numerically compatible with the table: 75.4 versus 35.0 tokens corresponds to a $2.154\times$ prompt or a 115.4\% increase.  However, the Abstract renders this comparison through malformed nested math, and the paper does not define whether 75.4 includes the system prompt, query, retrieved passage, and formatting.  Similarly, 14.45~s is called generation latency even though the claimed overhead is attributed partly to retrieval; retrieval, tokenization, prefill, and decoding times are not separated, and no repeated-run uncertainty is reported.

BM25 has only 19.20\% Recall@1.  Passage construction, Vietnamese tokenization, query formulation, BM25 implementation and parameters, relevance mapping, and retrieval-failure handling remain unspecified.  Oracle Gold-Context RAG is represented only by BERTScore 0.8002, without RefPref, prompt length, latency, or paired analysis.  The claim that retrieval noise causes the lower result is not established without retrieval-conditioned generation or an oracle comparison on all endpoints.  No paired CI/test supports the stated 8.60-point steering advantage over BM25 RAG, and ``zero-context-memory advantage'' is asserted without a memory measurement.

\findinglabel{Why it matters.} One under-specified, low-recall retriever cannot support a general comparison with RAG, and the systems conclusion mixes quantities with different boundaries.

\findinglabel{Suggested fix.} Publish the corpus construction and retrieval configuration; evaluate stronger lexical and dense retrievers; report category-wise Recall@$k$, results conditional on retrieval success, and complete oracle metrics.  Compare methods on identical questions with paired uncertainty.  Report total prompt tokens and component-wise timing/memory measurements with repeated-run intervals, and describe the result as one BM25 configuration.

\subsection*{8. Timing, ablation, and directional-control conclusions lack adequate uncertainty}

\classification{Statistical omission and overinterpretation}

\findinglabel{Location.} Experimental Setup; Table~\texttt{tab:baselines}; factorial ablation; Discussion; Conclusion.

\findinglabel{Problem.} Every non-RAG condition and all 12 ablation rows report exactly 7.32~s, yet the number of timing repetitions, warm-up policy, hardware sharing, prompt/output lengths, dispersion, and measurement boundaries are absent.  This does not establish ``no measurable'' overhead or equivalence.  Schedule differences are as small as 0.0001--0.0018 BERTScore, but no paired intervals or multiplicity-aware ablation analysis is given.  The paper acknowledges that these results are exploratory while still using them to support the method choice.

Positive steering changes RefPref by $+4.20$ points from baseline, whereas negative steering changes it by $-10.20$ points.  This asymmetry is not analyzed, so the sign control supports direction dependence but not a symmetric factuality axis.  The layer sweep likewise reports only a range and then infers distributed truthfulness and robustness without per-layer uncertainty.

\findinglabel{Why it matters.} Small schedule rankings may be noise, and the efficiency, symmetry, and layer-robustness claims are stronger than the reported evidence.

\findinglabel{Suggested fix.} Run repeated component-wise latency and peak-memory measurements with an equivalence margin.  Add paired ablation CIs and an exploratory multiplicity policy.  Release all layer-wise values and uncertainty, and analyze positive/negative asymmetry rather than calling it confirmation of a truthfulness axis.

\subsection*{9. Dataset, vector, metric, and hook documentation remains insufficient for reproduction}

\classification{Reproducibility limitation and likely leakage risk}

\findinglabel{Location.} benchmark construction; vector estimation; metric definition; Experimental Setup; implementation-facing hook expressions.

\findinglabel{Problem.} ``Expert-structured'' remains undefined, with no annotation protocol, counter-answer construction rules, clinical review fraction, agreement, adjudication, deduplication, source identifiers, or licensing.  The category counts 165, 160, and 175 sum to the 500-item evaluation subset but appear in the full 14,700-record construction paragraph without being identified as subset counts.  Question-disjoint splitting does not rule out drug/source-entry or error-template leakage, and the 500-item sampling rule and seed are absent.

The direction is extracted from the final token of a completed answer but injected into early autoregressive states without prompt-only, initial-token, pooled, or independent-error-template controls.  Rep-4 lacks a formula and tokenizer; Vietnamese ROUGE tokenization, complete BERTScore configuration, tie counts, statistical software calls, hook module/indexing, modified tensor positions, prompt-forward behavior, KV-cache handling, software patch versions, seeds, and GPU count remain unspecified.  Equation~4 maps ties to zero through the indicator, while the prose calls them ``neither''; their denominator and treatment in paired tests are not stated.

\findinglabel{Why it matters.} Leakage, tie handling, vector construction, and low-level hook choices can materially change every reported result and prevent independent replication.

\findinglabel{Suggested fix.} Add a dataset card and source/template-grouped split manifests, identify the category counts and evaluation sampling seed, document expert validation and licensing, test alternative vector extraction positions, and provide complete metric/statistical configurations.  Release executable hook pseudocode, exact environment details, per-question outputs, and all random seeds.

\section*{Equation, notation, sign, branch, and dimensional audit}

The direction $\mu_l^+-\mu_l^-$ and addition of $+\alpha(t)v_{\text{steer}}^{(l)}$ are sign-consistent with the intended positive intervention.  Unit normalization makes the vector dimensionally compatible with the hidden state when $\alpha(t)$ is an activation-scale coefficient.  The linear schedule and positive-$\alpha_0$ values of $D_{\text{continuous}}$, $D_{\text{cutoff}}$, and
\[
D_{\text{decay}}=\frac{K+1}{2}|\alpha_0|
\]
are arithmetically correct.  No inverse-trigonometric functions, coordinate transformations, or branch/quadrant choices occur.

The matched-dose definition is not sign-consistent over the stated negative-steering domain.  $D_{\text{decay}}$ is nonnegative, while
\[
\alpha_{\text{match}}=D_{\text{decay}}/T
=\frac{K+1}{2T}\alpha_0
\]
is negative when $\alpha_0<0$; consequently, the claim $\sum_t\alpha_{\text{match}}=D_{\text{decay}}$ cannot hold in that case.  Define $\alpha_{\text{match}}=\operatorname{sgn}(\alpha_0)D_{\text{decay}}/T$ and match $\sum_t|\alpha_{\text{match}}|$, or restrict this control to $\alpha_0>0$.

Additional concrete issues are the Layer~8/$l=16$ contradiction, pre-hook/post-hook drift between $h_l^{(t)}$ and $\hat h_l^{(t)}$, and overload of $N$ for training examples, category sizes, evaluation prompts, and random directions.  Decoding steps are discrete and should be written $t\in\{1,\ldots,100\}$ rather than $t\in[1,100]$.  Terminology drifts among ``truthfulness,'' ``factual preference,'' RefPref, and clinical accuracy, and among ``norm elevation,'' ``saturation,'' ``stabilization,'' and ``manifold departure.''

\section*{Meaningful editorial, compilation, and reference findings}

\subsection*{The manuscript does not complete a clean LaTeX build}

\findinglabel{Location.} Table~\texttt{tab:long\_gen}; caption of Table~\texttt{tab:mcnemar}; clean pdfLaTeX log.

\findinglabel{Problem.} The natural-completion float begins with \texttt{table} but ends with \texttt{table*}, causing a fatal environment-mismatch error under a halt-on-error build.  The significance-table caption contains an unescaped ampersand in ``Statistical Significance \& Pairwise Comparison Matrix,'' producing a second alignment error when compilation is allowed to continue.  The recovered build also reports the natural table approximately 61.44~pt too wide for its one-column placement.

\findinglabel{Why it matters.} A clean submission build produces no PDF, and error recovery can conceal broken float placement and table layout.

\findinglabel{Suggested fix.} Use a matching \texttt{table*} environment for the full-width natural table or redesign it for one column, escape the caption ampersand, and rebuild twice with halt-on-error.  Inspect the final log for overfull boxes and unresolved references.

\subsection*{The Abstract and experimental list contain malformed LaTeX/prose}

\findinglabel{Location.} Abstract; Experimental Setup.

\findinglabel{Problem.} The Abstract contains the nested math fragment beginning ``\texttt{\$context expansion (\$75.4\$ prompt tokens}'' and ending with a repeated ``prompt context expansion.''  The compiled output treats ordinary words as mathematical variables and produces severe spacing artifacts.  The Abstract is also a single overloaded sentence containing the complete RAG report.  The two evaluation modes in Experimental Setup are plain numbered paragraphs rather than a LaTeX list.

\findinglabel{Why it matters.} The headline summary is visibly malformed and difficult to interpret, while the experimental structure is not semantically typeset.

\findinglabel{Suggested fix.} Replace the fragment with plain prose such as ``75.4 versus 35.0 prompt tokens (a 115.4\% increase),'' shorten the RAG sentence, and use an \texttt{enumerate} environment for the evaluation modes.

\subsection*{Citation fit remains uneven}

\findinglabel{Location.} Introduction; bibliography entries \texttt{ref10}, \texttt{ref11}, and \texttt{ref12}.

\findinglabel{Problem.} \texttt{ref10} is the LoRA paper rather than direct evidence that SFT mitigates the stated medical hallucinations.  \texttt{ref11} studies robustness to irrelevant retrieved context but does not establish that internal activation pathways cause evidence ignoring.  \texttt{ref12} concerns catastrophic forgetting generally rather than the asserted medical-language-model failure mode.

\findinglabel{Why it matters.} Direct prior evidence, general background, and mechanistic hypotheses are presented with similar certainty.

\findinglabel{Suggested fix.} Add directly matched sources or qualify these statements as motivation and hypotheses.

\section*{Proofing sweep}

The bundled mechanical scan was run on both the source and the error-recovered PDF.  Its only hit was the semicolon before the proper name ``McNemar,'' which is not a capitalization error.  Manual equation-adjacent, table, Abstract, and bibliography checks identified these bounded defects:

\begin{itemize}
    \item \textbf{Activation Results:} replace the duplicated sign ``$++1.97$'' with ``$+1.97$'' and avoid ``within 3.01 units'' when the intended point is a 3.01-unit undershoot.
    \item \textbf{Natural table:} match the \texttt{table}/\texttt{table*} environment and fit the table to its actual width.
    \item \textbf{Significance-table caption:} escape the ampersand.
    \item \textbf{Abstract:} remove the malformed nested math and repeated ``prompt context expansion'' phrase.
    \item \textbf{Experimental Setup:} convert the two plain numbered paragraphs to an \texttt{enumerate} environment and replace the universal 100.00\% EOS statement with the 99.80\%--100.00\% condition-specific range.
\end{itemize}

The bibliography was checked for duplicated punctuation, malformed quotation punctuation, inconsistent capitalization, and obvious citation-format glitches; no additional mechanical defect was identified.

\section*{Items requiring external verification}

\begin{itemize}
    \item Verify all formulary-derived reference answers and synthetic counter-answers against the official source, current clinical guidance, qualified clinician review, and applicable reuse permissions.
    \item Validate multilingual BERTScore for Vietnamese clinical negation, quantities, units, contraindications, and interaction mechanisms before interpreting RefPref as factual alignment.
    \item Verify novelty against current token-dependent, scheduled, and dose-controlled activation-steering research.
    \item Verify the characterization of Qwen2.5-7B-Instruct as an SLM and the privacy/local-hospital deployment claims using explicit definitions, a threat model, and measured resource requirements.
    \item Verify the BM25 configuration and stronger Vietnamese retrieval baselines before generalizing the comparison with RAG.
\end{itemize}

\section*{Highest-priority revision sequence}

\begin{enumerate}
    \item Recover the actual placebo outputs and remove the impossible mean/range, $\sigma$, and randomization claims.
    \item Separate the natural and bounded paired analyses and regenerate all confidence intervals, McNemar tests, and multiplicity adjustments from their own outputs.
    \item Resolve Layer~8 versus $l=16$, specify pre/post-hook logging, and redesign the activation diagnostic around an operational mechanism on controlled token contexts.
    \item Repair the LaTeX float/caption errors and the malformed Abstract so the manuscript completes a clean build.
    \item Present natural completion as a metric trade-off and add paired uncertainty for RefPref, BERTScore, and Rep-4.
    \item Add clinician-adjudicated correctness/harm endpoints or narrow every clinical and hallucination claim to reference preference.
    \item Strengthen and fully specify RAG, timing, ablation, layer, dataset, leakage, metric, vector, and hook analyses.
\end{enumerate}

\end{document}